%%%%%%%%%%%%%%%%%%%%%%%%%%%%%%%%%%%%%%%%%%%%%%%%%%%%%%%
% Filename: "PARAMETERS.tex"
% -----------------------------------------------------
% To use the template, read the file "README.tex".
% -----------------------------------------------------
%
% Define all parameters in this file. They are used in
% various locations throughout the template.
%
%%%%%%%%%%%%%%%%%%%%%%%%%%%%%%%%%%%%%%%%%%%%%%%%%%%%%%%

% Author:
% - Use "~" to separate your names and your initials.
% - If you do not have a middle name, just use your
%   single name and initial (e.g., "A.")
\def\authorOfThesis{Dr. med. Andreas Roland Grüntzig}
\def\authorPrename{Andreas~Roland}
\def\authorSurname{Grüntzig}
\def\authorInitial{A.~R.}

% Institute:
\def\fakultaet{Medizinische Fakultät der Universität Zürich}
\def\institute{Department für Innere Medizin}
\def\institution{Universitätsspital Zürich}
\def\instituteAddress{Rämistrasse 100\\CH-8091 Zürich}

% Title page:
\def\director{Prof. Dr. med. Institutsdirektor}
\def\titleOfThesis{Die perkutane transluminale Rekanalisation chronischer Arterienverschlüsse}
\def\birthplace{Dresden, Deutschland}
\def\dateOfThesis{1. Juni 1977}

% Imprint:
% - You can enter up to five supervisors.
% - Leave blank if not needed.
\def\supervisorFirst{PD Dr. med. Erster Betreuer}
\def\supervisorSecond{Dr. med. Zweiter Betreuer}
\def\supervisorThird{}
\def\supervisorFourth{}
\def\supervisorFifth{}

% Your representative ("Fachvertreter"):
\def\representative{Prof. Dr. med. Fachvertreter}

% Copy shop where the thesis was printed:
% - Leave blank if not needed.
\def\printShop{Copy Shop\\Bahnhofplatz 1\\CH-8001 Zürich}
\def\printShopUrl{http://www.copyshop.ch/}

% Binding correction:
% - Space that is lost due to the binding (about 5--6 mm for glue
%   binding). You can ask your print shop for the correct value.
% - Recheck the layout of the compiled document after changing this
%   parameter!
\def\bindingCorrection{0mm}

% Show lists (true/fals) after the Table of Contents:
% - The list of original articles will always be shown,
%   as it is an integral part of the thesis.
\def\showListOfFigures{true}
\def\showListOfTables{false}
\def\showListOfAcronyms{true}
\def\showIndex{false}

% Style of citations:
% - See www.overleaf.com/learn/latex/Biblatex_bibliography_styles
%   for other citation styles. However, "nature" is recommended.
\def\citationStyle{nature}

% Highlight your own name in references:
% - Define whether your own name should be highlighted in the
%   references. In the list of original articles as well as in
%   the summary of articles, it will be highlighted anyway.
\def\highlightOwnName{false}

% Document status (draft/final):
% - "draft" can be used while preparing the document. This will
%   speed-up compilation and mark layout errors with black boxes.
\def\status{final}

% What do you want to compile (title/habil)?
% - "title" compiles the DIN A5 title page for the cover of your
%    thesis
% - "habil" compiles the whole manuscript (files 0–8)
\def\whatToCompile{habil}

% start compilation (DO NOT CHANGE!!!)
\input{template/parts/\whatToCompile}